% Readaptado de Semana 07 Clase 02 (2025-A): subespacios con motivacion geometrica.
\documentclass[compacto,10pt]{pubciencias-notas}
\usepackage{comandosEPN}
\usepackage{tikz}
\usetikzlibrary{babel}
\autor{Departamento de Formación Básica}
\materia{Álgebra Lineal}
\nota{Semana 06 Clase 01}
\fecha{Semestre 2026-B}
\logouno[1.1cm]{Logo_EPN.eps}
\logodos[1.7cm]{Logo_DFB.eps}
\definecolor{colordef}{cmyk}{0.81,0.62,0.00,0.22}
\raggedbottom

\begin{document}
\encabezado
\vspace*{-8mm}
\setcounter{section}{3}
\section{Espacios vectoriales}
\subsection{Motivación: la recta en $\R^2$}

En la semana anterior estudiamos espacios vectoriales con distintas clases
de objetos matemáticos. Ahora mantendremos las operaciones de un espacio conocido y
preguntaremos: \textbf{¿en qué subconjuntos podemos seguir sumando y
multiplicando por cualquier escalar sin salir del conjunto?}
En los ejemplos geométricos usaremos las operaciones usuales y escalares reales.

\begin{ejem}[Desplazarse en una dirección]\label{sem06cl01:recta}
    Un dispositivo se desplaza únicamente en la dirección $d=(1,2)$,
    en ambos sentidos y con cualquier magnitud. Sus desplazamientos forman
    la recta
    \[
        L=\{t(1,2):t\in\R\}=\{(x,y)\in\R^2:y=2x\}.
    \]
    Observemos que, en cada vector de $L$, la componente en $y$ es el
    doble de la componente en $x$. Por ejemplo, $u=(1,2)$ y $v=(1/2,1)$
    pertenecen a $L$. Su suma $u+v=(3/2,3)$ también pertenece a $L$.

    Sumar desplazamientos de $L$ o multiplicarlos por cualquier escalar
    real conserva la relación $y=2x$ y no nos hace salir de la recta.
    Después de introducir la definición, probaremos que estas propiedades
    permiten afirmar que $L$ es un \emph{subespacio de $\R^2$}.
\end{ejem}

\begin{center}
    \begin{tikzpicture}[x=1.25cm,y=0.65cm]
        \draw[->] (-1.5,0)--(2.3,0) node[right] {$x$};
        \draw[->] (0,-2.6)--(0,4.7) node[above] {$y$};
        \draw[colordef!65,thick] (-1.2,-2.4)--(2.1,4.2)
            node[above right] {$L:y=2x$};
        \draw[->,very thick,colordef] (0,0)--(1,2)
            node[midway,left] {$u$};
        \draw[->,thick,orange!80!black] (1,2)--(1.5,3)
            node[midway,right] {$v$};
        \draw[->,thick,teal!70!black] (0,0)--(-1,-2)
            node[midway,left] {$-u$};
        \fill (1.5,3) circle (1.5pt) node[left] {$u+v$};
        \fill (0,0) circle (1.5pt) node[below right] {$0$};
    \end{tikzpicture}
\end{center}


\subsection{Definición y criterio de subespacio}

\begin{defi}[Subespacio vectorial]\label{sem06cl01:definicion}
    Sea $V$ un espacio vectorial real y sea $W\subseteq V$ no vacío.
    Decimos que $W$ es un \textbf{subespacio vectorial de $V$} si $W$,
    con la misma suma y el mismo producto por escalar de $V$, es también
    un espacio vectorial real.

    Las operaciones en $W$ son las restricciones
    \[
        \begin{aligned}
            +\colon W\times W&\longrightarrow W,
            &(u,v)&\longmapsto u+v,\\
            \cdot\colon\R\times W&\longrightarrow W,
            &(\alpha,u)&\longmapsto\alpha u.
        \end{aligned}
    \]
    Para que tengan estos codominios, sus resultados deben pertenecer a $W$.
\end{defi}

\begin{teo}[Criterio de subespacio]\label{sem06cl01:criterio}
    Sea $V$ un espacio vectorial real. Un subconjunto $W\subseteq V$
    es un subespacio de $V$ si y solo si satisface:
    \begin{enumerate}
        \item $0_V\in W$.
        \item \textbf{Cerradura bajo la suma:} para todo $u,v\in W$,
              se tiene $u+v\in W$.
        \item \textbf{Cerradura bajo el producto por escalar:} para todo
              $u\in W$ y todo $\alpha\in\R$, se tiene $\alpha u\in W$.
    \end{enumerate}
\end{teo}

Retomemos la recta $L$ del ejemplo \ref{sem06cl01:recta} y apliquemos
el criterio. Primero, $(0,0)\in L$, pues $0=2\cdot0$.
Si $p=(s,2s)$ y $q=(t,2t)$ son dos vectores cualesquiera de $L$, entonces
\[
    p+q=(s+t,2s+2t)=(s+t,2(s+t))\in L.
\]
La componente en $y$ de la suma sigue siendo el doble de la componente
en $x$. Además, para cualquier $\alpha\in\R$,
\[
    \alpha p=(\alpha s,2\alpha s)=(\alpha s,2(\alpha s))\in L.
\]
Se cumplen las tres condiciones; por tanto, $L$ es un subespacio de $\R^2$.

\begin{obs}[Qué debemos comprobar]
    No se repiten los ocho axiomas de espacio vectorial: ya se cumplen en
    el espacio $V$, en este ejemplo, $\R^2$. Se comprueba que el subconjunto conserva los
    resultados de sus operaciones. Para aceptar un subespacio hay que
    razonar con vectores y escalares arbitrarios; para rechazarlo basta
    un contraejemplo a cualquiera de las tres condiciones.
\end{obs}

Los conjuntos $\{0_V\}$ y $V$ son siempre subespacios de $V$:
el primero solo contiene el vector cero y el segundo es el espacio completo.

\begin{ejem}[Una recta trasladada]\label{sem06cl01:trasladada}
    La recta $M=\{(x,y)\in\R^2:y=2x+1\}$ es paralela a $L$,
    pero no contiene al origen. Por eso no es un subespacio.
    Además, $u=(0,1)\in M$, mientras que
    \[
        0u=(0,0)\notin M,\qquad u+u=(0,2)\notin M.
    \]
\end{ejem}

\begin{advertencia}
    Contener al origen es necesario, pero no suficiente. También hay que
    comprobar las dos cerraduras.
\end{advertencia}

\subsection{Subespacios de matrices y de polinomios}

El mismo criterio se aplica a los espacios de la semana 05. Ahora los
vectores serán matrices o polinomios; en cada caso, debemos identificar
el vector cero y comprobar que las operaciones conservan la condición
que define al subconjunto.

\begin{ejem}[Matrices diagonales]\label{sem06cl01:matrices}
    En el espacio $\Mat[\R]{2}{2}$, consideremos el subconjunto
    \[
        D=\left\{\begin{pmatrix}a&0\\0&b\end{pmatrix}:a,b\in\R\right\}.
    \]
    Sus elementos son las matrices diagonales de tamaño $2\times2$.
    Usamos la suma de matrices y el producto por un escalar real usuales.
    \begin{enumerate}
        \item La matriz nula pertenece a $D$: basta tomar $a=b=0$.
        \item Si $A,B\in D$, sus entradas fuera de la diagonal son cero.
              Por tanto,
              \[
                  A+B=
                  \begin{pmatrix}a&0\\0&b\end{pmatrix}
                  +\begin{pmatrix}c&0\\0&d\end{pmatrix}
                  =\begin{pmatrix}a+c&0\\0&b+d\end{pmatrix}\in D.
              \]
        \item Para cualquier $\alpha\in\R$ y cualquier $A\in D$,
              \[
                  \alpha A=\begin{pmatrix}\alpha a&0\\0&\alpha b\end{pmatrix}
                  \in D.
              \]
    \end{enumerate}
    Así, $D$ es un subespacio de $\Mat[\R]{2}{2}$. La condición que
    conservan ambas operaciones es tener ceros fuera de la diagonal.
\end{ejem}

\begin{ejem}[Polinomios que se anulan en un punto]\label{sem06cl01:polinomios}
    En el espacio $\Pol_2$ de los polinomios reales de grado a lo sumo 2,
    consideremos
    \[
        W=\{p\in\Pol_2:p(1)=0\}.
    \]
    Por ejemplo, $t-1$ y $t^2-1$ pertenecen a $W$:
    sus gráficas pasan por el punto $(1,0)$. Para probar que $W$ es un
    subespacio, sean $p,q\in W$ dos polinomios arbitrarios de la forma
    \[
        p(t)=a_0+a_1t+a_2t^2,\qquad
        q(t)=b_0+b_1t+b_2t^2,
    \]
    con coeficientes reales. Como ambos pertenecen a $W$, cumplen
    \[
        a_0+a_1+a_2=0,\qquad b_0+b_1+b_2=0.
    \]
    \begin{enumerate}
        \item El polinomio cero pertenece a $W$, pues su valor en 1 es 0.
        \item La suma de $p$ y $q$ es
              \[
                  (p+q)(t)=(a_0+b_0)+(a_1+b_1)t+(a_2+b_2)t^2.
              \]
              Es un polinomio de grado a lo sumo 2 y, al evaluar en 1,
              \[
                  \begin{aligned}
                      (p+q)(1)&=(a_0+a_1+a_2)+(b_0+b_1+b_2)\\
                              &=0+0=0.
                  \end{aligned}
              \]
              Por tanto, $p+q\in W$.
        \item Para cualquier $\alpha\in\R$, el producto por escalar es
              \[
                  (\alpha p)(t)=\alpha a_0+\alpha a_1t+\alpha a_2t^2.
              \]
              También tiene grado a lo sumo 2 y cumple
              \[
                  (\alpha p)(1)=\alpha(a_0+a_1+a_2)=\alpha\cdot0=0.
              \]
              Por tanto, $\alpha p\in W$.
    \end{enumerate}
    Por tanto, $W$ es un subespacio de $\Pol_2$. Sumar o escalar estos
    polinomios conserva el paso de sus gráficas por $(1,0)$.

\end{ejem}

\section{Ejercicios}

\begin{ejer}[Dibujar, decidir y justificar]
    Para cada conjunto, haga un bosquejo y decida si es un subespacio del
    espacio indicado, con las operaciones usuales. Si lo es, verifique
    las tres condiciones del teorema \ref{sem06cl01:criterio} con vectores
    arbitrarios; si no, dé un contraejemplo e identifique qué condición falla.
    \begin{enumerate}
        \item $W_1=\{(x,y)\in\R^2:x+y=0\}$.
        \item $W_2=\{(x,y)\in\R^2:x\ge0,\ y\ge0\}$.
        \item $W_3=\{(x,y,z)\in\R^3:z=0\}$.
        \item $W_4=\{(x,y,z)\in\R^3:z=1\}$.
        \item $W_5=\{t(1,2):t\ge0\}$, en $\R^2$ (una semirrecta).
              Analice por separado si contiene al cero y si cumple cada
              una de las dos cerraduras.
    \end{enumerate}
\end{ejer}

\begin{ejer}[Velocidades permitidas en una plataforma]
    Una plataforma admite velocidades del conjunto
    \[
        W=\{(v_x,v_y,v_z)\in\R^3:v_x+v_y-v_z=0\}.
    \]
    \begin{enumerate}
        \item Interprete geométricamente $W$ y compruebe que
              $a=(1,0,1)$ y $b=(0,1,1)$ pertenecen a él.
        \item Calcule $2a-b$ y compruebe que sigue siendo admisible.
              Explique qué permite afirmar esto sobre las dos cerraduras
              y por qué un solo cálculo no basta para demostrarlas en general.
        \item Si se añade la restricción $v_z\ge0$, ¿el conjunto de
              velocidades sigue siendo un subespacio? Dé un contraejemplo.
    \end{enumerate}
\end{ejer}



\begin{ejer}[Salida individual]
    Un compañero afirma: ``Si un conjunto contiene al cero y todos los
    múltiplos de cada uno de sus vectores, entonces es un subespacio''.
    Explique qué condición falta y construya un contraejemplo en $\R^2$.
\end{ejer}

\begin{ejer}
    Con las operaciones usuales de la semana 05, decida si son subespacios:
    \begin{enumerate}
        \item $U=\{A\in\Mat[\R]{2}{2}:a_{12}=a_{21}\}$, en
              $\Mat[\R]{2}{2}$.
        \item $T=\{p\in\Pol_2:p(0)=0\}$, en $\Pol_2$.
        \item $S=\{p\in\Pol_2:p(0)=1\}$, en $\Pol_2$.
    \end{enumerate}
    Justifique con el criterio de subespacio, sin volver a verificar los
    ocho axiomas de espacio vectorial.
\end{ejer}

\end{document}
